% ============================================================
%  Rapport de projet IA | Onitama
%  Université Paris Cité
% ============================================================
\documentclass[12pt, a4paper]{article}

% --- Encodage & langue ---
\usepackage[utf8]{inputenc}
\usepackage[T1]{fontenc}
\usepackage[french]{babel}
\usepackage{csquotes}

% --- Mise en page ---
\usepackage[top=2.5cm, bottom=2.5cm, left=2.8cm, right=2.8cm]{geometry}
\usepackage{setspace}
\onehalfspacing

% --- Typographie ---
\usepackage{lmodern}
\usepackage{microtype}

% --- Couleurs & graphiques ---
\usepackage[dvipsnames]{xcolor}
\usepackage{graphicx}
\usepackage{float}
\usepackage{tikz}
\usetikzlibrary{shapes, arrows.meta, positioning, calc}

% --- Tableaux ---
\usepackage{booktabs}
\usepackage{array}
\usepackage{multirow}
\usepackage{tabularx}
\usepackage{colortbl}

% --- Code source ---
\usepackage{listings}
\lstset{
  language=Python,
  basicstyle=\ttfamily\small,
  keywordstyle=\color{MidnightBlue}\bfseries,
  commentstyle=\color{gray}\itshape,
  stringstyle=\color{Maroon},
  numberstyle=\tiny\color{gray},
  numbers=left,
  numbersep=6pt,
  breaklines=true,
  frame=single,
  rulecolor=\color{gray!40},
  backgroundcolor=\color{gray!5},
  showstringspaces=false,
  tabsize=4,
  captionpos=b,
  literate=
    {é}{{\'{e}}}1 {è}{{\`{e}}}1 {ê}{{\^{e}}}1 {ë}{{\"e}}1
    {à}{{\`{a}}}1 {â}{{\^{a}}}1 {ù}{{\`{u}}}1 {û}{{\^{u}}}1
    {î}{{\^{i}}}1 {ï}{{\"i}}1 {ô}{{\^{o}}}1 {ç}{{\c{c}}}1
    {À}{{\`{A}}}1 {É}{{\'{E}}}1 {È}{{\`{E}}}1 {Ê}{{\^{E}}}1,
}

% --- Maths ---
\usepackage{amsmath}
\usepackage{amssymb}

% --- Algorithmes ---
\usepackage[ruled, vlined, linesnumbered, french]{algorithm2e}

% --- Liens hypertexte ---
\usepackage[colorlinks=true, linkcolor=MidnightBlue, citecolor=MidnightBlue, urlcolor=MidnightBlue]{hyperref}

% --- En-têtes / pieds de page ---
\usepackage{fancyhdr}
\pagestyle{fancy}
\fancyhf{}
\rhead{\small Projet IA -- Onitama}
\lhead{\small Université Paris Cité}
\cfoot{\thepage}

% --- Utilitaires ---
\usepackage{enumitem}
\usepackage{caption}
\usepackage{subcaption}

% ============================================================
%  Commandes personnalisées
% ============================================================
\newcommand{\todo}[1]{\colorbox{yellow!60}{\textbf{[À COMPLÉTER : #1]}}}
\definecolor{bleuonitama}{RGB}{30, 100, 200}
\definecolor{rougeonitama}{RGB}{200, 40, 40}

% --- Nouvelles extensions conseillées ---
\usepackage{url}
\usepackage{fontawesome5}
\usepackage[default]{sourcesanspro}
\usepackage[most]{tcolorbox}
\usepackage{tocloft}

% --- Réglages sommaire ---
\setlength{\cftbeforesecskip}{2pt}
\setlength{\cftbeforesubsecskip}{0pt}
\setcounter{tocdepth}{2}

% --- Couleurs additionnelles pour tcolorbox ---
\definecolor{main}{HTML}{1F4E79}
\definecolor{accent}{HTML}{2E86AB}
\definecolor{light}{HTML}{F5F7FA}

% ============================================================
\begin{document}

% ============================================================
%  PAGE DE TITRE
% ============================================================
\begin{titlepage}
\centering

\vspace*{1cm}

% --- Université ---
{\large\textsc{\textcolor{main}{\faUniversity\ Université Paris Cité}}} \\
\vspace{0.2cm}
{\normalsize \faGraduationCap\ Licence Informatique -- UE Intelligence Artificielle}

\vspace{1cm}

\rule{0.7\linewidth}{1.2pt}

\vspace{1.5cm}

% --- Titre ---
{\Huge\bfseries\color{main} IA et Onitama}

\vspace{0.5cm}

% --- Sous-titre ---
{\Large\itshape Exploration des algorithmes Minimax et de l'élagage $\alpha\text{-}\beta$}

\vspace{1.5cm}

\rule{0.7\linewidth}{1.2pt}

\vspace{2cm}
\begin{minipage}{0.9\textwidth}
    \centering
    \large
    % Étudiant 1 à gauche
    \begin{minipage}[t]{0.45\textwidth}
        \centering
        \textcolor{main}{\faUser} \\
        \textbf{Simeon TAKENNE MEKEM} \\
        \textit{\small Informatique et Application \& Mathématiques et Applications}
    \end{minipage}
    \hfill % Pousse le deuxième bloc vers la droite
    % Étudiant 2 à droite
    \begin{minipage}[t]{0.45\textwidth}
        \centering
        \textcolor{main}{\faUser} \\
        \textbf{Amine ZEROUALI} \\
        \textit{\small Informatique et Application \& Mathématiques et Applications}
    \end{minipage}
\end{minipage}
\vspace{2cm}

% --- Enseignante ---
{\normalsize 
\faChalkboardTeacher\ \textbf{Responsable de l'UE :} Elise Bonzon
}

\vspace{0.5cm}

% --- GitHub ---
\faGithub\ 
\href{https://github.com/Amine-Zerouali/onitama-minimax-explorer}
{\color{accent}Code source du projet}

\vfill

% --- Résumé ---
\begin{minipage}{0.85\textwidth}
\small
\centering
\textbf{\textcolor{main}{Résumé :}}
Ce document présente la conception d’un moteur de jeu pour \textit{Onitama} ainsi que l'implémentation de trois niveaux de difficulté basés sur l'algorithme Minimax. L’efficacité de l’élagage $\alpha\text{-}\beta$ est démontrée par une analyse du nombre de nœuds explorés et une étude statistique portant sur 1000 parties par configuration. Les résultats confirment la supériorité des heuristiques positionnelles sur le simple comptage matériel.
\end{minipage}

\vspace{1cm}

{\small Février -- Mai 2026}

\end{titlepage}

% ============================================================
%  TABLE DES MATIÈRES
% ============================================================
\newpage
\begingroup
\setstretch{1.0}
\tableofcontents
\endgroup
\newpage
\newpage

% ============================================================
\section{Introduction}
% ============================================================

Ce rapport présente notre projet d'intelligence artificielle appliqué au jeu de société
\textit{Onitama}. L'objectif est de concevoir un moteur de jeu complet ainsi que plusieurs agents
artificiels de niveaux de difficulté croissants, en s'appuyant sur les algorithmes classiques
de la théorie des jeux.

Dans ce rapport, nous décrirons dans un premier temps les règles du jeu et les raisons de
notre choix (Section~\ref{sec:jeu}), nous présenterons ensuite les détails d'implémentation du
moteur de jeu (Section~\ref{sec:moteur}), puis les algorithmes d'IA et nos trois heuristiques
(Section~\ref{sec:ia}). Nous analyserons enfin les résultats de nos tournois (Section~\ref{sec:tournoi})
avant de conclure (Section~\ref{sec:conclusion}).

% ============================================================
\section{Présentation du jeu}
\label{sec:jeu}
% ============================================================

\subsection{Description générale}

\textit{Onitama} est un jeu de stratégie à deux joueurs, déterministe et à information parfaite,
créé par Shimpei Sato. Il se joue sur un plateau de $5 \times 5$ cases. Chaque joueur dispose
de cinq pièces : quatre \textbf{élèves} et un \textbf{maître}.

\begin{figure}[H]
  \centering
  \begin{tikzpicture}[scale=0.9]

    % --- OMBRE GLOBALE ---
    \begin{scope}[shift={(0.08,-0.08)}]
      \fill[black!10, rounded corners=0.1cm] (0,0) rectangle (5,5);
    \end{scope}

    % --- FOND DU PLATEAU ---
    \fill[gray!10, rounded corners=0.1cm] (0,0) rectangle (5,5);

    % --- GRILLE ---
    \draw[gray!60, line width=0.8pt] (0,0) grid (5,5);
    \draw[black!70, thick, rounded corners=0.1cm] (0,0) rectangle (5,5);

    % --- TEMPLES ---
    % Bleu (haut)
    \fill[bleuonitama!25] (2,4) rectangle (3,5);
    \draw[bleuonitama!80, thick] (2,4) rectangle (3,5);

    % Rouge (bas)
    \fill[rougeonitama!25] (2,0) rectangle (3,1);
    \draw[rougeonitama!80, thick] (2,0) rectangle (3,1);

    % --- JOUEUR BLEU ---
    \foreach \x in {0,1,3,4} {
      % Ombre
      \fill[black!20] (\x+0.55,4.45) circle (0.35);
      % Pion
      \fill[bleuonitama] (\x+0.5,4.5) circle (0.35);
      \draw[black!50] (\x+0.5,4.5) circle (0.35);
    }

    % Maître bleu
    \fill[black!25] (2.55,4.45) circle (0.45);
    \fill[bleuonitama!90] (2.5,4.5) circle (0.45);
    \draw[black, thick] (2.5,4.5) circle (0.45);
    \node[white, font=\bfseries\small] at (2.5,4.5) {M};

    % --- JOUEUR ROUGE ---
    \foreach \x in {0,1,3,4} {
      % Ombre
      \fill[black!20] (\x+0.55,0.45) circle (0.35);
      % Pion
      \fill[rougeonitama] (\x+0.5,0.5) circle (0.35);
      \draw[black!50] (\x+0.5,0.5) circle (0.35);
    }

    % Maître rouge
    \fill[black!25] (2.55,0.45) circle (0.45);
    \fill[rougeonitama!90] (2.5,0.5) circle (0.45);
    \draw[black, thick] (2.5,0.5) circle (0.45);
    \node[white, font=\bfseries\small] at (2.5,0.5) {M};

    % --- AXES ---
    \foreach \x in {0,...,4} {
      \node[font=\scriptsize, gray] at (\x+0.5, 5.25) {\x};
    }

    \foreach \y [count=\i from 0] in {4.5, 3.5, 2.5, 1.5, 0.5} {
      \node[font=\scriptsize, gray] at (-0.25, \y) {\i};
    }

  \end{tikzpicture}
  \caption{État initial du plateau. \newline Chaque joueur possède quatre pions et un maître (noté M). Les cases centrales colorées représentent les temples. \newline Convention \texttt{plateau[y][x]} : la ligne $y=0$ correspond au haut du plateau (Joueur Bleu) et $y=4$ au bas (Joueur Rouge).}
  \label{fig:plateau_init}
\end{figure}


\subsection{Règles du jeu}

Au début de chaque partie, cinq cartes de mouvement sont tirées aléatoirement parmi seize.
Deux cartes sont attribuées à chaque joueur et la dernière, appelée \textbf{carte extra}, déterminera quant à elle
quel joueur commence :
\begin{itemize}
    \item si sa couleur est bleue, le Joueur 1 commence
    \item si elle est rouge, c'est le Joueur 2.
\end{itemize}


\paragraph{Déroulement d'un tour.} Le joueur actif choisit l'une de ses deux cartes et
déplace l'une de ses pièces selon les mouvements autorisés par cette carte. Les déplacements
sont relatifs à l'orientation du joueur (le Joueur~2 joue dans le sens symétrique). La carte jouée prend la place de la carte extra. Cette dernière est récupérée par le joueur, tandis que la carte qu'il vient de jouer sera intégrée à la main de l'adversaire à la fin de son prochain tour.

\paragraph{Conditions de victoire.} Un joueur remporte la partie selon l'un des deux chemins :
\begin{itemize}[leftmargin=*]
  \item \textbf{Voie de la Pierre} : capturer le maître adverse ;
  \item \textbf{Voie du Courant} : amener son propre maître sur le temple adverse
        (case $(2, 0)$ pour le Joueur~1, case $(2, 4)$ pour le Joueur~2).
\end{itemize}

\subsection{Intérêt stratégique et justification du choix}

Onitama présente plusieurs propriétés qui en font un choix particulièrement adapté au cadre
de ce projet :

\begin{itemize}
  \item \textbf{Déterminisme et information parfaite} : toutes les informations sont visibles
        des deux joueurs, ce qui satisfait les hypothèses de l'algorithme Minimax.
  \item \textbf{Espace d'états suffisamment grand} : la combinaison du plateau $5\times5$ et
        du tirage aléatoire des cinq cartes parmi seize génère un nombre de parties très élevé,
        rendant toute résolution par force brute impossible.
        \item \textbf{Originalité et concision} : les règles sont simples à programmer mais le facteur de branchement reste significatif (pouvant atteindre jusqu'à une quarantaine de coups possibles par tour selon les cartes disponibles et la position des pièces), garantissant une profondeur de recherche intéressante.
        \item \textbf{Double condition de victoire} : la coexistence des deux voies de victoire
        crée des situations stratégiques riches, propices à l'élaboration d'heuristiques
        nuancées.
\end{itemize}

% ============================================================
\section{Implémentation du moteur de jeu}
\label{sec:moteur}
% ============================================================

\subsection{Représentation de l'état}

L'état du jeu est encapsulé dans la classe \texttt{Onitama}. Le plateau est représenté par une
matrice Python $5 \times 5$ d'entiers signés :

\begin{center}
\begin{tabular}{cl}
  \toprule
  Valeur & Signification \\
  \midrule
  \texttt{0}  & Case vide \\
  \texttt{1}  & Élève du Joueur~1 (Bleu) \\
  \texttt{2}  & Maître du Joueur~1 (Bleu) \\
  \texttt{-1} & Élève du Joueur~2 (Rouge) \\
  \texttt{-2} & Maître du Joueur~2 (Rouge) \\
  \bottomrule
\end{tabular}
\end{center}

Ce codage signé permet de déterminer l'appartenance d'une pièce à un joueur par un simple
test du signe : \texttt{piece * joueur > 0}.

Les cartes sont stockées dans le dictionnaire global \texttt{CARTES}, dont chaque entrée
contient la liste des déplacements relatifs (paires $(\Delta x, \Delta y)$) et la couleur
associée. Les mains des deux joueurs sont représentées par deux listes de noms de cartes, et
la carte extra par un unique identifiant. La variable \texttt{tour} ($\pm 1$) indique quel
joueur doit jouer.

\subsection{Génération des coups}

La méthode \texttt{obtenir\_coups\_possibles(joueur)} parcourt le plateau pour identifier les
pièces du joueur actif, puis, pour chaque pièce et pour chaque carte disponible, calcule les
positions cibles en appliquant le facteur directionnel (\texttt{sens = ±1}). Un coup est valide
si la case cible est dans les bornes du plateau et n'est pas occupée par une pièce alliée.

\begin{lstlisting}[caption={Extrait de la génération de coups légaux.}]
for (x, y) in mouvements:
    nx = j + (x * sens)
    ny = i + (y * sens)
    if 0 <= nx < 5 and 0 <= ny < 5:
        cible = self.plateau[ny][nx]
        if cible * joueur <= 0:   # case vide ou piece adverse
            coups.append({"carte": carte, "de": (j, i), "vers": (nx, ny)})
\end{lstlisting}

\paragraph{Note sur les règles officielles :} Une légère simplification a été apportée concernant 
le cas très rare où un joueur se retrouve sans aucun mouvement légal (totalement bloqué par 
ses propres pièces et les bords du plateau). Selon les règles strictes d'Onitama, le joueur 
doit alors passer son tour mais défausser une carte pour l'échanger avec la carte extra. 
Dans notre implémentation, pour maintenir l'efficacité de l'arbre de recherche, le tour passe 
simplement à l'adversaire sans rotation des cartes. Ce cas limite étant exceptionnellement 
rare en pratique, son impact sur l'évaluation globale des heuristiques est considéré comme négligeable.

% ============================================================

\subsection{Application et annulation des coups}

La méthode \texttt{jouer\_coup} déplace la pièce, effectue l'échange de cartes et inverse
le tour courant. Elle retourne la pièce capturée et l'ancienne carte extra afin de permettre
l'annulation du coup par \texttt{annuler\_coup}, qui restaure l'état complet en ordre inverse.
Ce mécanisme est indispensable à l'exploration récursive de l'arbre de
jeu sans copie coûteuse de l'état.

\subsection{Détection des états terminaux}

La méthode \texttt{verifier\_victoire} parcourt le plateau et retourne \texttt{1}, \texttt{-1}
ou \texttt{None} selon que le Joueur~1 a gagné, le Joueur~2 a gagné ou que la partie est
encore en cours. Les deux conditions de victoire sont vérifiées simultanément :

\begin{lstlisting}[caption={Détection des conditions de victoire.}]
if piece == 2:           # Maitre bleu
    master_p1_vivant = True
    if i == 4 and j == 2:    # Voie du Courant (temple rouge)
        return 1
elif piece == -2:        # Maitre rouge
    master_p2_vivant = True
    if i == 0 and j == 2:    # Voie du Courant (temple bleu)
        return -1
if not master_p1_vivant: return -1   # Voie de la Pierre
if not master_p2_vivant: return  1
\end{lstlisting}

Bien que le moteur n'impose pas de limite intrinsèque, la boucle principale du jeu (lancer\_partie) impose une limite de 1000 tours pour éviter les parties théoriquement infinies, déclarant alors la partie nulle.

\subsection{Tests unitaires}
\label{sec:tests}

Afin de garantir la fiabilité du moteur de jeu avant d'y greffer l'algorithme Minimax,
plusieurs tests unitaires ont été implémentés à l'aide de la bibliothèque standard \texttt{unittest}
de Python. 

Ces tests ciblent les composants critiques du jeu et s'assurent qu'aucune régression
n'est introduite lors de l'optimisation du code. Les points clés validés sont :

\begin{itemize}
  \item \textbf{L'initialisation de l'état} : vérification de la bonne distribution des cinq
        cartes (deux par joueur, une au centre) et du bon placement des pièces.
  \item \textbf{La détection des états terminaux} : simulation de positions spécifiques 
        (comme le placement forcé d'un maître sur le temple adverse) pour s'assurer que
        la méthode \texttt{verifier\_victoire} retourne le bon gagnant.
  \item \textbf{La réversibilité des coups (Make/Unmake)} : un test crucial vérifie que l'enchaînement
        des méthodes \texttt{jouer\_coup} et \texttt{annuler\_coup} restaure exactement l'état
        du plateau, les mains des joueurs, la carte extra et le tour courant. C'est cette validation qui 
        garantit que l'arbre de recherche Minimax ne corrompt pas l'état du jeu lors de son parcours.
\end{itemize}

L'exécution de cette suite de tests valide le socle logique du projet et permet aux heuristiques
de s'appuyer sur des états de jeu rigoureusement exacts.

% ============================================================
\section{Intelligence artificielle}
\label{sec:ia}
% ============================================================

\subsection{Algorithme Minimax}

L'algorithme Minimax est un algorithme de recherche dans l'arbre de jeu adapté aux jeux à
somme nulle à deux joueurs. Il s'appuie sur la récursivité :

\begin{itemize}
  \item Si l'état est terminal ou si la profondeur maximale est atteinte, retourner la valeur
        de l'heuristique ;
  \item Si c'est le tour du joueur MAX (Bleu), retourner la valeur maximale parmi les enfants ;
  \item Si c'est le tour du joueur MIN (Rouge), retourner la valeur minimale parmi les enfants.
\end{itemize}



Dans notre implémentation, les deux fonctions récursives \texttt{max\_value} et
\texttt{min\_value} partagent les paramètres $\alpha$ et $\beta$ dès leur premier appel depuis
\texttt{minimax\_decision}, ce qui intègre nativement l'élagage $\alpha\text{-}\beta$ (voir
Section~\ref{sec:alphabeta}).

\begin{algorithm}[H]
\DontPrintSemicolon
\SetAlgoLined
\caption{Minimax (version MAX)}
\KwIn{état $s$, profondeur $d$, $\alpha$, $\beta$, heuristique $h$}
\KwOut{valeur de l'état}
\If{$s$ est terminal}{
  \Return{score terminal}
}
\If{$d = 0$}{
  \Return{$h(s)$}
}
$v \leftarrow -\infty$\;
\ForEach{coup $a$ applicable dans $s$}{
  $s' \leftarrow \text{jouer}(s, a)$\;
  $v \leftarrow \max\{v,\;\textsc{MinValue}(s', d-1, \alpha, \beta, h)\}$\;
  $\text{annuler}(s, a)$\;
  \If{$v \geq \beta$}{\Return{$v$} \tcp{élagage}}
  $\alpha \leftarrow \max\{\alpha, v\}$\;
}
\Return{$v$}
\end{algorithm}

\subsection{Élagage $\alpha\text{-}\beta$}
\label{sec:alphabeta}

L'élagage $\alpha\text{-}\beta$ maintient deux bornes globales :
\begin{itemize}[leftmargin=*]
  \item $\alpha$ : meilleure valeur garantie pour MAX (initialisé à $-\infty$) ;
  \item $\beta$ : meilleure valeur garantie pour MIN (initialisé à $+\infty$).
\end{itemize}



Lorsque $v \geq \beta$ dans un nœud MAX (respectivement $v \leq \alpha$ dans un nœud MIN),
les branches restantes sont élagués car elles ne peuvent améliorer le résultat. Dans le
meilleur des cas, cet élagage réduit la complexité temporelle de $O(b^m)$ à $O(b^{m/2})$,
doublant ainsi la profondeur de recherche atteignable.

Notre implémentation mélange aléatoirement les coups (\texttt{random.shuffle}) avant
l'exploration, ce qui introduit de la variété dans les parties tout en empêchant un biais
systématique lié à l'ordre de génération.

\begin{lstlisting}[caption={Élagage $\alpha\text{-}\beta$ dans \texttt{max\_value}.}]
v = float('-inf')
for coup in coups:
    piece_capturee, ancienne_extra = self.jouer_coup(coup)
    v = max(v, self.min_value(profondeur - 1, alpha, beta, heuristique))
    self.annuler_coup(coup, piece_capturee, ancienne_extra)
    if v >= beta:          # Elagage beta
        return v
    alpha = max(alpha, v)
return v
\end{lstlisting}

\subsection{Fonctions d'évaluation (heuristiques)}
\label{sec:heuristiques}

En cas de non-atteinte d'un état terminal dans les limites de profondeur, l'algorithme recourt
à une fonction d'évaluation pour estimer la valeur d'une position. Nous avons implémenté trois
heuristiques, toutes exprimées du point de vue du Joueur~1 (MAX) :

\medskip
\begin{center}
\begin{tabular}{clp{7cm}}
  \toprule
  Identifiant & Nom court & Description \\
  \midrule
  H1 & Matériel basique & Comptage pondéré fixe des pièces \\
  H2 & Matériel dynamique & Valeur marginale décroissante des pions \\
  H3 & Avancée (H2 + centre) & H2 enrichie d'un bonus positionnel \\
  \bottomrule
\end{tabular}
\end{center}

\subsubsection{H1}

Cette heuristique attribue un score fixe à chaque pièce encore en jeu :

\begin{equation}
  \text{H1}(s) = 10 \times (\#\text{élèves bleus}) + 100 \times (\#\text{maîtres bleus})
                 - 10 \times (\#\text{élèves rouges}) - 100 \times (\#\text{maîtres rouges})
\end{equation}

Elle est intentionnellement simple : chaque pion vaut autant, qu'il soit le dernier ou le
quatrième d'une équipe. Cette heuristique correspond au niveau de difficulté \textbf{Facile}.

\subsubsection{H2}

H2 introduit une \textit{valeur marginale décroissante} : plus les pions se raréfient, plus
chacun prend de la valeur. Ce schéma correspond à la réalité stratégique d'Onitama, où perdre
son dernier élève est une catastrophe bien plus grave que de perdre un pion parmi quatre.

\begin{equation}
  \text{H2}(s) = f(\#\text{élèves bleus}) - f(\#\text{élèves rouges})
\end{equation}

où $f$ est définie par :

\begin{equation}
  f(n) = \begin{cases}
    0  & n = 0 \\
    8  & n = 1 \\
    14 & n = 2 \\
    18 & n = 3 \\
    20 & n \geq 4
  \end{cases}
\end{equation}

Cette heuristique constitue le niveau \textbf{Moyen}.

\subsubsection{H3}

H3 enrichit H2 d'un bonus de positionnement qui récompense le contrôle du centre du plateau :

\begin{equation}
  \text{H3}(s) = \text{H2}(s)
    + \sum_{\text{pièce } p} \text{signe}(p) \times \left(2 \cdot \mathbb{1}_{p \in \text{zone centrale}} + 3 \cdot \mathbb{1}_{p \in \text{case centre}}\right)
\end{equation}
où $\mathbb{1}_A$ désigne la fonction indicatrice de l'ensemble $A$.

La \textit{zone centrale} est définie comme $\{(x,y) \mid 1 \leq x \leq 3,\, 1 \leq y \leq 3\}$
et la \textit{case centre} comme $(2,2)$.

Ce bonus traduit l'importance tactique du contrôle du centre, qui accroît la mobilité des
pièces et rapproche physiquement le maître des deux temples adverses. Cette heuristique
constitue le niveau \textbf{Difficile}.

\subsection{Profils des IA et niveaux de difficulté}

Les trois niveaux de difficulté sont définis par la combinaison d'une profondeur de recherche
et d'une heuristique :

\begin{center}
\begin{tabular}{clcc}
  \toprule
  Niveau & Nom & Profondeur & Heuristique \\
  \midrule
  1 & Facile   & 3 & H1 \\
  2 & Moyen    & 4 & H2 \\
  3 & Difficile & 5 & H3 \\
  \bottomrule
\end{tabular}
\end{center}

\noindent\textit{Remarque :} pour les expériences comparant uniquement les heuristiques (équité
du test), une profondeur commune de~4 est utilisée pour les trois profils.

% ============================================================
\section{Analyse expérimentale}
\label{sec:analyse_exp}
% ============================================================

\subsection{Temps de calcul par profondeur}
\label{sec:temps}

\begin{table}[H]
  \centering
  \caption{Temps moyen de calcul d'un coup (en secondes) selon la profondeur et l'heuristique.}
  \label{tab:temps}
  \begin{tabular}{crrr}
    \toprule
    Profondeur & H1 (s) & H2 (s) & H3 (s) \\
    \midrule
    1 & 0.0001 (±0.0000) & 0.0001 (±0.0000) & 0.0001 (±0.0000) \\
    2 & 0.0002 (±0.0001) & 0.0002 (±0.0001) & 0.0005 (±0.0003) \\
    3 & 0.0012 (±0.0005) & 0.0014 (±0.0006) & 0.0029 (±0.0013) \\
    4 & 0.0069 (±0.0046) & 0.0053 (±0.0048) & 0.0166 (±0.0096) \\
    5 & 0.0452 (±0.0369) & 0.0475 (±0.0300) & 0.1449 (±0.0905) \\
    6 & 0.1688 (±0.0846) & 0.1538 (±0.1131) & 0.5727 (±0.4783) \\
    7 & 1.0830 (±0.7311) & 2.0322 (±2.3179) & 2.8081 (±1.8955) \\
    \bottomrule
  \end{tabular}
\end{table}

On observe que le temps de calcul croît de façon quasi-exponentielle avec la profondeur,
conformément à la complexité théorique $O(b^{m/2})$ après élagage $\alpha\text{-}\beta$.
H3 est systématiquement plus lente que H1 et H2 en raison du parcours supplémentaire
du plateau pour calculer le bonus positionnel. À profondeur 5, H3 prend environ 3 fois
plus de temps que H1 ou H2 ; à profondeur 7, la différence s'accentue (2,8~s contre 1,1~s
pour H1). En pratique, la profondeur 4 représente le meilleur compromis performance/temps
pour H3 (16~ms en moyenne), tandis que H1 et H2 restent exploitables jusqu'à la profondeur 6.

\begin{figure}[H]
  \centering
  \includegraphics[width=0.9\linewidth]{Image/temps_calcul.png}
  \caption{Évolution du temps de calcul en fonction de la profondeur pour les trois heuristiques
           (échelle logarithmique sur l'axe des ordonnées).}
  \label{fig:temps}
\end{figure}


\subsection{Nombre de nœuds explorés}
\label{sec:noeuds}

Le tableau~\ref{tab:noeuds} présente le nombre moyen de nœuds visités par l'algorithme
Minimax avec et sans élagage $\alpha\text{-}\beta$, mesuré avec l'heuristique H3 à profondeurs croissantes.
Ces mesures ont été obtenues en ajoutant un compteur global incrémenté à chaque appel
à \texttt{max\_value} et \texttt{min\_value}.

\begin{table}[H]
  \centering
  \caption{Nombre moyen de nœuds explorés avec et sans élagage $\alpha\text{-}\beta$ (H3).}
  \label{tab:noeuds}
  \begin{tabular}{lrrr}
    \toprule
    Profondeur & Sans élagage & Avec élagage $\alpha\text{-}\beta$ & Ratio\\
    \midrule
    2 &    175 &    66 &  2{,}64 \\
    3 &  2\,441 &   580 &  4{,}21 \\
    4 & 42\,591 & 3\,109 & 13{,}70 \\
    5 & 852\,038 & 24\,545 & 34{,}71 \\
    \bottomrule
  \end{tabular}
\end{table}

\begin{figure}[H]
  \centering
  \includegraphics[width=0.9\linewidth]{Image/noeuds_explores.png}
  \caption{Impact de l'élagage $\alpha\text{-}\beta$ sur le nombre de nœuds explorés selon la profondeur
           (H3, échelle logarithmique).}
  \label{fig:noeuds}
\end{figure}

\subsection{Impact de l'élagage $\alpha\text{-}\beta$ sur l'efficacité}

L'élagage $\alpha\text{-}\beta$ réduit théoriquement la complexité temporelle de $O(b^m)$ à
$O(b^{m/2})$ dans le meilleur cas. Dans notre contexte, le facteur de branchement moyen
se situe autour d'une dizaine de coups par tour (bien qu'il puisse théoriquement atteindre la quarantaine dans de rares cas).

Les résultats du tableau~\ref{tab:noeuds} et de la figure~\ref{fig:noeuds} confirment
l'efficacité croissante de l'élagage avec la profondeur : le ratio Sans/Avec passe de
$\times2{,}6$ à profondeur~2 à $\times34{,}7$ à profondeur~5. Cette progression super-linéaire
s'explique par l'effet cumulatif de l'élagage sur plusieurs niveaux de l'arbre : plus l'arbre
est profond, plus la propagation des bornes $\alpha$ et $\beta$ permet d'éliminer de branches
entières.

Le mélange aléatoire des coups (\texttt{random.shuffle}) a un effet ambigu sur l'efficacité
de l'élagage. Dans le meilleur cas théorique ($O(b^{m/2})$), les coups sont présentés du
meilleur au pire, ce qui maximise les coupures. Notre ordre aléatoire s'écarte de cet optimum
et induit une variabilité importante des ratios d'une position à l'autre ; en contrepartie, il
garantit la diversité des parties et évite que l'IA explore toujours les mêmes variantes en
premier.

% ============================================================
\section{Tournoi entre les intelligences artificielles}
\label{sec:tournoi}
% ============================================================

\subsection{Protocole expérimental}

Pour évaluer les performances relatives de nos trois niveaux de difficulté, nous avons organisé
un tournoi automatisé à l'aide du script \texttt{generer\_graphiques.py}. Chaque paire d'IA
s'affronte dans $1000$ parties. Le taux de victoire du Joueur Bleu (J1) est enregistré pour
chaque croisement.

Les résultats sont présentés sous la forme d'une matrice de chaleur $3 \times 3$ : chaque
cellule $(i, j)$ indique le pourcentage de victoires de l'IA~$i$ (jouant Bleu) contre
l'IA~$j$ (jouant Rouge).

\subsection{Résultats}

La figure~\ref{fig:matrice} illustre les résultats finaux du tournoi. Les taux de victoire présentés correspondent au pourcentage de parties remportées par l'IA jouant les pièces bleues (Joueur 1) face à l'IA jouant les pièces rouges (Joueur 2).

\begin{figure}[H]
  \centering
  \includegraphics[width=0.65\textwidth]{Image/matrice_resultats.png}
  \caption{Matrice de chaleur des taux de victoire du Joueur Bleu (\%).}
  \label{fig:matrice}
\end{figure}

\subsection{Évolution du taux de victoire}

La figure~\ref{fig:evolution} présente l'évolution du taux de victoire cumulé du joueur Bleu pour trois affrontements différents entre les IA (Facile, Moyenne et Difficile), sur 1000 parties. L'axe des abscisses est représenté en échelle logarithmique afin de mettre en évidence la convergence rapide des taux.

\begin{figure}[H]
  \centering
  \includegraphics[width=0.75\textwidth]{Image/evolution_victoires.png}
  \caption{Évolution du taux de victoire cumulé du joueur Bleu selon les paires d'IA, sur 1000 parties (axe $x$ en échelle logarithmique). Les taux se stabilisent respectivement à 89,7\%, 99,1\% et 96,7\%.}
  \label{fig:evolution}
\end{figure}

Dès les premières dizaines de parties, on observe quelques oscillations (particulièrement visibles sur la courbe orange) avant que les taux ne se stabilisent progressivement. À l'issue des 1000 parties, les courbes convergent vers leurs valeurs finales, démontrant une hiérarchie claire entre les niveaux :

\begin{itemize}
    \item \textbf{Difficile vs Facile (courbe rouge) :} Le taux atteint \textbf{99,1\%}, illustrant une domination quasi totale de l'IA Difficile.
    \item \textbf{Difficile vs Moyen (courbe bleue) :} Le taux se stabilise à \textbf{96,7\%}, confirmant l'apport significatif de l'heuristique H3 et de la profondeur 5 par rapport au niveau Moyen.
    \item \textbf{Moyen vs Facile (courbe orange) :} Le taux converge vers \textbf{89,7\%}, prouvant la nette supériorité de la profondeur 4 (H2) sur l'IA de base.
\end{itemize}

Les lignes pointillées colorées matérialisent ces taux finaux pour chaque affrontement, tandis que la ligne grise en tirets marque le seuil de 50\% (égalité parfaite).

\subsection{Analyse des résultats}

\subsubsection{Hiérarchie des niveaux}

Les résultats confirment parfaitement la hiérarchie attendue entre les trois niveaux de
difficulté (Facile $<$ Moyen $<$ Difficile) :

\begin{itemize}
  \item \textbf{Moyen bat Facile} avec un taux de 89{,}9\% (Bleu), l'avantage de H2 sur H1
        et la profondeur supérieure (4 vs 3) produisent une très forte domination.
  \item \textbf{Difficile bat Moyen} avec un taux de 98{,}0\% (Bleu), le bonus positionnel H3
        combiné à la profondeur~5 confère un avantage quasi absolu.
  \item \textbf{Difficile bat Facile} avec un taux de 98{,}4\% (Bleu), ce qui illustre 
        l'écart massif de niveau de jeu entre les deux extrêmes.
\end{itemize}

Les matchs miroir (même IA contre elle-même) donnent des taux remarquablement équilibrés et proches de 50\% :
48{,}1\% pour Facile vs Facile, 49{,}4\% pour Moyen vs Moyen et 49{,}0\% pour Difficile
vs Difficile. L'asymétrie potentielle liée au fait de jouer en premier ou en second semble 
minime, confirmant l'équité de l'affrontement lorsque les deux joueurs ont le même niveau.

\subsubsection{Robustesse statistique}

Avec 1000 parties par croisement, l'intervalle de confiance à 95\% sur un taux de victoire $p$
est approximativement :

\begin{equation}
  \hat{p} \pm 1{,}96 \sqrt{\frac{\hat{p}(1-\hat{p})}{1000}}
\end{equation}

Pour $\hat{p} = 0{,}5$, cet intervalle vaut environ $\pm 3{,}1$ points de pourcentage, bien plus
serré qu'avec 50 parties ($\pm 13{,}9$ pp). Les résultats des matchs miroir (compris entre 48{,}1\% et 49{,}4\%) s'inscrivent parfaitement dans cet intervalle de $50\% \pm 3{,}1\%$, validant l'hypothèse d'une égalité parfaite. En revanche,
les résultats extrêmes ($\geq 89\%$ ou $\leq 10\%$) sont statistiquement incontestables et ne
requièrent pas davantage de tirages pour conclure à une domination claire.

\subsubsection{Influence des cartes aléatoires}

Le tirage aléatoire des cinq cartes introduit une source de variabilité initiale dans chaque partie.
Certaines combinaisons de cartes favorisent nettement un style de jeu (par exemple, des cartes
offrant de grands sauts favorisent une stratégie d'encerclement), ce qui peut biaiser ponctuellement
le résultat d'une manche isolée.

Cette variabilité est d'ailleurs visible dans les fortes oscillations initiales sur la courbe d'évolution du
taux de victoire (figure~\ref{fig:evolution}). Cependant, les résultats des matchs miroirs, très proches de 50\%, démontrent que les 1000 parties par croisement sont suffisantes pour échantillonner une
très grande variété de tirages de cartes. La loi des grands nombres a opéré, lissant l'impact de la distribution initiale des cartes et permettant d'isoler l'impact pur des heuristiques et de la profondeur d'exploration.

% ============================================================
\section{Utilisation de l'IA générative}
\label{sec:ia_generative}
% ============================================================

Dans le cadre de ce projet, nous avons eu recours à plusieurs outils d'intelligence artificielle 
générative, avec une répartition stricte des tâches, pour nous assister dans les phases de 
développement et de rédaction : les modèles \textbf{Claude Sonnet 4.6} et \textbf{Gemini 3.1 Pro} (via leurs 
interfaces web), ainsi que l'éditeur de code \textbf{Cursor}.

\begin{itemize}
  \item \textbf{Type d'utilisation :} 
        L'éditeur \textbf{Cursor} a été utilisé spécifiquement pour la documentation, facilitant la 
        rédaction et la structuration des commentaires directement dans le code source. 
        Le modèle \textbf{Claude} a été sollicité pour la partie programmation, principalement pour 
        les scripts annexes (tests unitaires, génération de graphiques) et pour comprendre comment 
        implémenter efficacement le parallélisme en Python (\texttt{concurrent.futures}). 
        Enfin, \textbf{Gemini} a été principalement dédié au rapport : vérification mathématique des 
        données de benchmark, relecture et mise en forme globale en \LaTeX.
        
  \item \textbf{Rôle effectif :} 
        Chaque IA a joué un rôle délimité. Cursor a permis de commenter le code de manière propre et lisible. 
        Claude a agi comme un tuteur technique, débloquant les difficultés liées à la gestion du 
        \textit{multiprocessing} nécessaire pour accélérer les milliers de parties du tournoi. 
        Gemini a endossé le rôle de relecteur et d'assistant d'édition, permettant d'aligner rigoureusement 
        les analyses textuelles avec les sorties brutes de la console (calcul des ratios, vérification des 
        intervalles de confiance) et de restructurer le document.
        
  \item \textbf{Analyse critique (apports et limites) :} 
        La segmentation des tâches entre ces outils s'est révélée très pertinente. L'apport principal 
        réside dans le gain de temps sur des tâches chronophages (commentaires, formatage des tableaux \LaTeX) 
        et dans l'aide à la compréhension de l'architecture parallèle (avec Claude). En revanche, 
        l'utilisation de ces LLM exige une rigueur stricte dans les instructions (le \textit{prompting}) : 
        lors de la rédaction, il a été indispensable d'imposer à l'IA de ne formuler aucune conclusion sans 
        recevoir au préalable les journaux d'exécution (\textit{logs}) exacts, afin d'éviter toute 
        hallucination de résultats ou de statistiques.
\end{itemize}


% ============================================================
\section{Difficultés rencontrées}
\label{sec:difficultes}
% ============================================================

Le développement de ce moteur de jeu et de son intelligence artificielle a soulevé plusieurs 
défis techniques et conceptuels, que nous avons résolus par diverses optimisations :

\begin{itemize}
  \item \textbf{Symétrie directionnelle des cartes :} Contrairement aux échecs où les pièces ont des 
        règles de mouvement absolues, les cartes d'Onitama sont relatives au joueur. Nous avons résolu 
        cette difficulté en introduisant un multiplicateur de direction (\texttt{sens = 1} pour Bleu, 
        \texttt{-1} pour Rouge) appliqué systématiquement lors de la génération des coordonnées cibles $(nx, ny)$.
  
  \item \textbf{Implémentation du Make/Unmake :} Pour que l'algorithme Minimax explore des millions 
        de nœuds rapidement, il était impossible de copier le plateau entier à chaque coup (\texttt{deepcopy}). 
        La difficulté principale résidait dans l'échange de la \textit{carte extra} : il a fallu s'assurer 
        que la méthode \texttt{annuler\_coup} réinsère la bonne carte dans la bonne main, retire l'ancienne 
        extra, et replace la pièce capturée éventuelle. Les tests unitaires ont été indispensables pour 
        valider cette mécanique.

  \item \textbf{Temps d'exécution à haute profondeur :} Le langage Python étant interprété, le parcours 
        de l'arbre devient très lent au-delà de la profondeur 5, particulièrement avec l'heuristique H3 
        qui nécessite de parcourir le plateau pour calculer les bonus de position. Pour générer nos 
        benchmarks (comme le tournoi de 1000 parties), le temps d'exécution devenait prohibitif. Nous avons 
        pallié ce problème en utilisant le module \texttt{concurrent.futures.ProcessPoolExecutor}, permettant 
        de paralléliser les parties sur tous les cœurs du processeur.

  \item \textbf{Gestion des cas limites et parties nulles :} Deux cas limites majeurs ont été traités. 
        D'une part, pour éviter que l'IA ne s'enferme dans des cycles infinis, une limite stricte de 
        1000 tours a été imposée dans la boucle de jeu principale. D'autre part, le cas très rare où un 
        joueur n'a aucun coup légal a été simplifié : pour maintenir l'efficacité de l'arbre Minimax, 
        le moteur passe directement le tour sans obliger l'IA à modéliser la défausse d'une carte.
        
  \item \textbf{Lissage du hasard des cartes :} Lors de nos premiers tests du tournoi sur 50 ou 250 parties, 
        les taux de victoires montraient une forte variance liée au tirage initial des cartes. Ce n'est qu'en 
        augmentant l'échantillon à 1000 parties par croisement que la loi des grands nombres a permis 
        d'isoler la valeur réelle de nos heuristiques des "bons tirages" ponctuels.
\end{itemize}

% ============================================================
\newpage
\section{Conclusion}
\label{sec:conclusion}
% ============================================================

Ce projet nous a permis de mettre en pratique les concepts fondamentaux de l'intelligence
artificielle appliquée aux jeux de réflexion, au travers du jeu \textit{Onitama}.

Nous avons conçu un moteur de jeu complet reposant sur un mécanisme efficace de \textit{make/unmake},
permettant d'explorer l'arbre de jeu sans copier l'état à chaque nœud. L'algorithme Minimax
avec élagage $\alpha\text{-}\beta$ a été intégré dès la racine, réduisant substantiellement le nombre
de nœuds évalués.

Trois heuristiques de complexité croissante ont été développées : H1 (comptage matériel fixe),
H2 (valeur marginale décroissante) et H3 (H2 augmentée d'un bonus positionnel centré).
Les tournois automatisés ont confirmé la hiérarchie attendue entre les niveaux de difficulté :
Facile $<$ Moyen $<$ Difficile, avec des écarts de performance très marqués (allant de 
près de 90\% à plus de 98\% de victoires d'un niveau sur l'inférieur). L'élagage $\alpha\text{-}\beta$ s'est 
révélé particulièrement efficace, réduisant de près d'un facteur 35 ($\times 34{,}7$) le nombre de nœuds 
explorés à profondeur~5. H3 offre les meilleures performances stratégiques au prix d'un surcoût 
de calcul d'un facteur 3 environ par rapport à H1 et H2, coût qui reste acceptable jusqu'à la 
profondeur~5 (environ 145~ms par coup).

Plusieurs pistes d'amélioration sont envisageables pour repousser les limites de profondeur. 
Pour l'approche Minimax, l'ordonnancement des coups avant l'élagage (tri par score heuristique), 
l'utilisation de tables de transposition pour la mémoïsation des états~\cite{memoisation}, ou encore 
l'exploration itérative en profondeur (IDA*) permettraient d'optimiser drastiquement les 
coupures $\alpha\text{-}\beta$~\cite{IDA}. Alternativement, l'implémentation d'une approche 
probabiliste comme la recherche arborescente Monte-Carlo (MCTS)~\cite{mcts}, 
particulièrement efficace sur des jeux à fort facteur de branchement, constituerait un 
prolongement naturel à ce projet.

% ============================================================
%  RÉFÉRENCES
% ============================================================
\newpage
\begin{thebibliography}{9}

\bibitem{onitama_rules}
Arcane Wonders.
\textit{Onitama | Livret de règles officiel}.
\newline Disponible sur : \url{https://www.arcanewonders.com/wp-content/uploads/2021/05/Onitama-Rulebook.pdf}

\bibitem{onitama_wiki}
Wikipédia.
\textit{Onitama} (Article en anglais).
\newline Disponible sur : \url{https://en.wikipedia.org/wiki/Onitama}

\bibitem{onitama_github}
Jack Adamson.
\textit{Onitama} (Dépôt GitHub).
\newline Disponible sur : \url{https://github.com/jackadamson/onitama}

\bibitem{memoisation}
C\#ESS Enginer.
\textit{Episode\_08\_11: La table de transposition} (Vidéo en ligne).
YouTube, 2021.
\newline Disponible sur : \url{https://www.youtube.com/watch?v=gBOIFDDUjcE}

\bibitem{IDA}
Mahesh Huddar.
\textit{
Iterative Deepening A Star algorithm IDA Star | IDA* – Artificial intelligence by Mahesh Huddar} (Vidéo en ligne).
YouTube, 2024.
\newline Disponible sur : \url{https://www.youtube.com/watch?v=BUHc8p5Mpdo}

\bibitem{mcts}
John Levine.
\textit{Monte Carlo Tree Search} (Vidéo en ligne).
YouTube, 2017.
\newline Disponible sur : \url{https://www.youtube.com/watch?v=UXW2yZndl7U}

\end{thebibliography}

% ============================================================
%  ANNEXES
% ============================================================
\appendix
\newpage
\section{Code source principal}
\label{ann:code}

\subsection{Évaluation heuristique H2  - Matériel dynamique}

\begin{lstlisting}[caption={Heuristique H2 : valeur marginale décroissante des pions.}]
def evaluation_materiel_dynamique(self):
    def points_par_pions(compte):
        if compte == 0: return 0
        if compte == 1: return 8
        if compte == 2: return 14
        if compte == 3: return 18
        if compte >= 4: return 20

    pions_bleus = pions_rouges = 0
    maitre_bleu = maitre_rouge = False

    for y in range(5):
        for x in range(5):
            p = self.plateau[y][x]
            if p == 1:  pions_bleus  += 1
            elif p == 2:  maitre_bleu   = True
            elif p == -1: pions_rouges += 1
            elif p == -2: maitre_rouge  = True

    if not maitre_bleu:  return -10000
    if not maitre_rouge: return  10000
    return points_par_pions(pions_bleus) - points_par_pions(pions_rouges)
\end{lstlisting}

\subsection{Script de génération des graphiques}

Le script \texttt{generer\_graphiques.py} automatise les benchmarks et les tournois, 
puis produit quatre graphiques :
\begin{itemize}
  \item \texttt{temps\_calcul.png} : temps moyen par coup selon la profondeur ;
  \item \texttt{noeuds\_explores.png} : nœuds explorés avec et sans élagage $\alpha\text{-}\beta$ ;
  \item \texttt{matrice\_resultats.png} : matrice de chaleur $3\times3$ des taux de victoire ;
  \item \texttt{evolution\_victoires.png} : courbe du taux de victoire cumulé selon le nombre de parties(axe $x$ en échelle logarithmique).
\end{itemize}

\noindent Pour l'exécuter : \texttt{python3 generer\_graphiques.py}

\section{Dictionnaire des 16 cartes}
\label{ann:cartes}

\begin{table}[H]
  \centering
  \caption{Les 16 cartes d'Onitama avec leurs déplacements $(\Delta x, \Delta y)$ et leur couleur de départ.}
  \label{tab:cartes}
  \begin{tabular}{llp{7.5cm}}
    \toprule
    \textbf{Carte} & \textbf{Couleur} & \textbf{Déplacements $(\Delta x, \Delta y)$} \\
    \midrule
    Tiger   & Bleu  & $(0,2)$, $(0,-1)$ \\
    Dragon  & Rouge & $(-2,1)$, $(-1,-1)$, $(1,-1)$, $(2,1)$ \\
    Frog    & Rouge & $(-2,0)$, $(-1,1)$, $(1,-1)$ \\
    Rabbit  & Bleu  & $(-1,-1)$, $(1,1)$, $(2,0)$ \\
    Crab    & Bleu  & $(-2,0)$, $(0,1)$, $(2,0)$ \\
    Elephant & Rouge & $(-1,1)$, $(-1,0)$, $(1,0)$, $(1,1)$ \\
    Goose   & Bleu  & $(-1,1)$, $(-1,0)$, $(1,0)$, $(1,-1)$ \\
    Rooster & Rouge & $(-1,-1)$, $(-1,0)$, $(1,0)$, $(1,1)$ \\
    Monkey  & Bleu  & $(-1,1)$, $(-1,-1)$, $(1,1)$, $(1,-1)$ \\
    Mantis  & Rouge & $(-1,1)$, $(0,-1)$, $(1,1)$ \\
    Horse   & Rouge & $(-1,0)$, $(0,1)$, $(0,-1)$ \\
    Ox      & Bleu  & $(0,1)$, $(0,-1)$, $(1,0)$ \\
    Crane   & Bleu  & $(0,1)$, $(-1,-1)$, $(1,-1)$ \\
    Boar    & Rouge & $(0,1)$, $(-1,0)$, $(1,0)$ \\
    Eel     & Bleu  & $(-1,1)$, $(-1,-1)$, $(1,0)$ \\
    Cobra   & Rouge & $(-1,0)$, $(1,1)$, $(1,-1)$ \\
    \bottomrule
  \end{tabular}
\end{table}

\noindent\textit{Note :} les déplacements sont relatifs à la direction du joueur. Le Joueur~2
(Rouge) joue avec $\text{sens} = -1$, ce qui inverse les coordonnées $y$ par rapport au
Joueur~1. La couleur de la carte extra détermine quel joueur commence la partie.

\end{document}